\section{Dimension bounds for material derivatives}
\label{sec:divergence}

The Gaussian interpolation has supplied spatial derivatives of every
order using only the original $C^2$ potential.  We now need derivatives
along its transporting trajectories.  These are the derivatives that
control the error when a path integral is replaced by polynomial
quadrature.  Applying a Euclidean norm after each differentiation would
lose a new power of $d$ each time.  We will keep at least two tensor
slots free throughout the calculation and convert to a vector norm
only at the last step.

Retain the primitive parameters $L=\lambda r$ and $l=rL\le1/2$,
and the interpolation
$\xi_\rho=\mathcal L(\rho Z+\sqrt a\,N)$, where
$a=1-\rho^2$ and $0\le\rho<1$.  Expectations and $L^p$ norms in
this section use $\xi_\rho$ at the displayed probability time.  Along
a solution of $Y'_\rho=-rM_\rho(Y_\rho)$, the chain rule gives
\[
 \frac d{d\rho}Q(\rho,Y_\rho)
   =(\mathcal DQ)(\rho,Y_\rho),\qquad
 \mathcal D=\partial_\rho-rM_\rho\cdot\nabla.
\]
Thus the required path derivatives are iterates of $\mathcal D$.

The useful form of their identities involves the adjoint of a spatial
derivative.  If a tensor field $V$ has a distinguished coordinate slot
$k$, define its \emph{adjoint divergence on that slot} by
\[
 (\nabla^*V)_I
 =\sum_{k=1}^d\{-\partial_kV_{I,k}
                    +(\partial_kH_\rho)V_{I,k}\}.
\]
Here $I$ lists the slots that remain; the distinguished slot need not
have been last before a permutation.  For a compactly supported smooth
test function $\phi$, the defining integration-by-parts relation is
\[
 \int(\nabla^*V)_I\phi\,d\xi_\rho
   =\int\sum_{k=1}^dV_{I,k}\partial_k\phi\,d\xi_\rho.
\]
The next identities express material derivatives using this operation
and contractions of the normalized spatial jets
$U_q=\rho^{-q}\nabla^qM_\rho$ from
Proposition~\ref{prop:gaussian-jets}.

\subsection{Material derivatives of the conditional jets}

\begin{lemma}
\label{lem:material-identities}
The mean and its jets satisfy
\begin{align}
 \mathcal D M&=\nabla^*U_1,
 \label{eq:material-mean-identity}\\
 (\mathcal D U_q)_{i,I}
 &= (\nabla^*U_{q+1})_{i,I}
   +2r\rho\sum_{\varnothing\ne S\subseteq I}
       (U_{|S|})_{k,S}
       (U_{q-|S|+1})_{i,k,I\setminus S},\qquad q\ge1.
 \label{eq:material-jet-identity}
\end{align}
In the second identity $I$ is a labeled list of $q$ derivative slots;
subsets refer to positions in that list and the index $k$ is summed.
For any smooth tensor field $V$,
\begin{equation}
 \mathcal D(\nabla^*V)
 =\nabla^*(\mathcal DV)
   +r\nabla^*((DM_\rho)V),
 \label{eq:material-adjoint-identity}
\end{equation}
where the matrix acts on the slot removed by the divergence.  These
identities hold also at $\rho=0$ by continuous extension.
\end{lemma}

\begin{proof}
Differentiate
$\E[f_i(Z)\phi(Y)]=\int (M_\rho)_i\phi\,d\xi_\rho$.
The right side differentiates using the transporting velocity from
Proposition~\ref{prop:transport}.  On the left the random velocity is
$(Z-\rho Y)/a$.  Subtract its conditional mean from the flux multiplying
$\nabla\phi$.  The remaining coefficient is
$\Cov(f_i(Z),Z_k\mid Y)/a=(U_1)_{i,k}$ by
\eqref{eq:jet-cumulant}.  Integration by parts proves
\eqref{eq:material-mean-identity}.

For the higher identity, coordinate differentiation gives
\[
 \mathcal D\partial_I M_i
 =\partial_I\mathcal DM_i
 +r\sum_{\varnothing\ne S\subseteq I}
       (\partial_SM_k)\partial_k\partial_{I\setminus S}M_i.
\]
Insert \eqref{eq:material-mean-identity} and multiply by $\rho^{-q}$.
When all derivatives pass inside the adjoint divergence, the result
is $\nabla^*U_{q+1}$.  In the other terms they hit the score
$y+r\rho M_\rho$.  A derivative hitting its linear part contributes
$qU_q/\rho$, canceled by differentiating the prefactor $\rho^{-q}$.
The nonlinear score terms and the material/spatial commutator terms
are identical after scaling: each is
\[
 r\rho\sum_{\varnothing\ne S\subseteq I}
       (U_{|S|})_{k,S}(U_{q-|S|+1})_{i,k,I\setminus S}.
\]
This proves \eqref{eq:material-jet-identity}.  Each contraction retains
at least one free slot in each tensor.

Finally differentiate the adjoint identity against a transported
scalar test, componentwise in the other slots.  The terms involving
$\mathcal D\phi$ cancel.  The gradient commutator is
\[
 \mathcal D\partial_k\phi-\partial_k\mathcal D\phi
 =r(\partial_kM_j)\partial_j\phi.
\]
Moving it to the other side by the adjoint relation gives
\eqref{eq:material-adjoint-identity}.  All coefficients in the final
formulas are nonsingular at zero, so the extensions follow from
Proposition~\ref{prop:gaussian-jets}.
\end{proof}

\subsection{Contractions and the expressions they generate}

Before iterating the material identities, we isolate the norm estimate
that makes their product terms harmless.  Each coordinate slot represents
a copy of $\R^d$.  Contracting one slot of each of two tensors means
summing the product over one common coordinate, leaving every other
slot free.  The restriction to one slot is what lets us interpret the
operation as matrix composition.

\begin{lemma}[A single contraction]
\label{lem:split-contraction}
Let $A$ and $B$ be tensors, and contract one slot of $A$ with one slot
of $B$, leaving at least one free slot in each.  For any nontrivial
split of the resulting free slots there are nontrivial splits of $A$
and $B$ such that
\[
 |A*B|_{\mathrm{split}}
 \le |A|_{\mathrm{split}}|B|_{\mathrm{split}}.
\]
\end{lemma}

\begin{proof}
Call the two sides of the required output split the input and output
sides.  Orient the contracted index from one factor to the other.
If all free slots of one factor lie on the same side, its contracted
slot must be placed on the opposite side of that factor's unfolding.
These requirements are compatible: incompatibility would put every
free slot of the result on one side, contrary to the hypothesis.
If a factor already has free slots on both sides, it imposes no
requirement on the orientation.  With the resulting orientation, the
contraction is composition of the two matrix actions, with identities
on the unused input and output coordinates and with coordinate
permutations.  Identities and these permutations have operator norm
one.  Submultiplicativity proves the assertion.
\end{proof}

We can now specify exactly which tensor formulas will be estimated.
This also records which contractions may occur: contracting two slots
of a single tensor is not one of the allowed operations.

\begin{definition}[Tensor expressions and their weights]
\label{def:divergence-expressions}
A tensor expression is a finite formula built by the following rules.
\begin{enumerate}
 \item For each integer $q\ge1$, $U_q/L$ is an expression with $q+1$
       free slots and weight $q-1$.
 \item A permutation of the free slots preserves the weight.
 \item Two expressions may be contracted on one chosen slot each,
       provided each factor retains at least one free slot.  The
       resulting weight is the sum of their weights.
 \item An adjoint divergence on a chosen free slot is allowed when
       at least two free slots remain.  It increases the weight by one.
\end{enumerate}
The formula specifies its leaves, their orders, and the slots used at
each operation.  This finite syntactic data is independent of
$d,\rho,x,r$ and of the potential.  Repeated occurrences are treated
as separate factors, even when their values agree.
\end{definition}

For example, $\nabla^*(U_2/L)$ has two free slots and weight two,
whereas $\nabla^*(U_1/L)$ is not an expression under these rules:
it has vector output.  The latter is precisely the operation we will
reserve for the final $\sqrt d$ estimate.  If $E$ is an expression,
$\nabla^jE$ has $j$ additional free derivative slots, and a split may
place any of these slots on either side.  The next theorem bounds all
such splits simultaneously.  Its dependence on $j$ and $p$ is explicit
because its proof must later take a moment order growing with $\log d$.

For a related Gaussian result,
\citet[Theorem~1.1 and Corollary~1.4]{chen2026divergence} prove sharp
divergence inequalities with Hilbert--Schmidt derivative norms. Here we
require estimates under the generally non-Gaussian law $\xi_\rho$ that
retain every matrix split through the contractions and adjoint divergences
in Definition~\ref{def:divergence-expressions}.

\Needspace{12\baselineskip}
\begin{theorem}[Tensor divergence bound]
\label{thm:tensor-divergence}
Let $E$ be a tensor expression of weight $k$ in
Definition~\ref{def:divergence-expressions}.  There is a constant
$C_E\ge1$, depending only on that finite formula, such that for
every integer $j\ge0$, every real $2\le p<\infty$,
every $0\le\rho<1$, and every nontrivial split,
\begin{equation}
 \bigl\|\,|\nabla^jE|_{\mathrm{split}}\,\bigr\|_{L^p(\xi_\rho)}
 \le C_E^{j+1}
       \bigl(p+j+\log(d+1)\bigr)^{C_E(j+1)}
       (\sqrt a)^{-(k+j)}.
 \label{eq:network}
\end{equation}
The same $C_E$ works for all $j,p,d,\rho,x,r$.  In particular it is
not restricted to fixed derivative or moment orders.
\end{theorem}

\begin{proof}
We induct on the finite construction of $E$, proving all derivative
and moment orders simultaneously.  Permutations and single contractions
follow directly from the preceding lemma.  For a divergence, we first
integrate by parts in a high trace moment.  The resulting coordinate
sums can be arranged as a composition of matrix unfoldings with only
one closing trace.  We then choose the moment large enough to absorb
that trace's dimension factor.

For a normalized jet,
$\nabla^jU_q=\rho^jU_{q+j}$, so
Proposition~\ref{prop:gaussian-jets} proves the assertion after
enlarging a constant depending on $q$.  Permutations cause no change.
For a contraction, distribute the derivatives by Leibniz, use
Lemma~\ref{lem:split-contraction} on each summand, and use the
inductive estimates at exponent $2p$ on both factors.  The number of
Leibniz summands is at most $2^j$.  Their weights and derivative
orders add, giving \eqref{eq:network} with an enlarged fixed constant.

It remains to consider $E=\nabla^*V$, where $V$ has weight $k_V$ and
$E$ has at least two slots.  Commute the $j$ derivatives past the
divergence.  Apart from
\[
 T=\nabla^*(\nabla^jV),
\]
the terms are single contractions of $\nabla^{j-h}V$ with an
order-$h$ derivative of the score, $1\le h\le j$.  Both factors
retain slots.  The score derivative is bounded pointwise by a
constant of the form $C^h h^{Ch}(\sqrt a)^{-(h-1)}$, by
\eqref{eq:interpolation-score}--\eqref{eq:score-derivatives}.
Lemma~\ref{lem:split-contraction}, the inductive hypothesis, and
Leibniz multiplicities therefore give the required bound for all
these terms, even with exponent $k_V+j-1$ in place of the larger
allowed exponent $k_V+j+1$.  It remains to bound $T$.

\paragraph{Removing the adjoints.}
Fix an unfolding of $T$, with both sides nonempty, and an integer
$s\ge1$.  Pointwise,
\[
 \|T\|_{\op}^{2s}\le\tr((TT^T)^s).
\]
Expand the trace by coordinate indices.  It contains $2s$ labeled
primary factors, each an adjoint divergence of $\nabla^jV$.
Their free slots are joined in the alternating trace cycle; the two
bundles incident to each primary factor are nonempty.

Remove adjoint divergences one at a time by integration by parts.
For a chosen primary, its adjoint differentiates the product of the
other factors, never its own inside field.  If a derivative hits a
remaining adjoint, use
\[
 \partial_i(\nabla^*W)
 =\nabla^*(\partial_iW)
    +(\partial_i\partial_kH_\rho)W_k.
\]
The first term leaves that adjoint to be processed later.  The second
consumes it and creates a Hessian factor joining the two divergence
indices.  A later derivative may also hit a previously created
Hessian factor.  After all adjoints have been removed, a typical term
has primary factors $\nabla^{j+e_i}V$, $1\le i\le2s$, and $n_H$
factors $\nabla^{h_v}\nabla^2H_\rho$.  The bookkeeping identity is
\begin{equation}
 \sum_{i=1}^{2s}e_i+\sum_{v=1}^{n_H}h_v+2n_H=2s.
 \label{eq:network-degree-count}
\end{equation}
Indeed, a removed adjoint either adds one derivative or pairs with
another adjoint to create a Hessian.  Each step has at most $Cs$
possible target factors and at most two commutator terms, so there
are at most $(Cs)^{Cs}$ resulting terms.

We use a graph only to describe these explicit coordinate sums.
Its primary vertices are the $2s$ labeled tensor factors; a Hessian
vertex represents one factor $\nabla^{h_v}\nabla^2H_\rho$.
An edge represents a pair of tensor slots summed over a common
coordinate, and a bundle represents several such edges.
Three properties will control its contraction.  All original trace
bundles remain.  Every new direct edge connects distinct primary
factors, because an adjoint never differentiates its own inside
field.  Finally, each Hessian factor has at least two distinct
primary neighbors, the pair that created it, and there are no
Hessian-to-Hessian edges.  Later derivatives only add further
primary-to-Hessian edges.  These statements remain true if some
derivatives had already been commuted inside an unprocessed adjoint.

\paragraph{Opening the trace.}
Choose one bundle in the original trace cycle and cut it, leaving
its $b\ge1$ indices unsummed at the two ends.  All remaining original
bundles form a chain through the $2s$ primary factors.  Order these
factors along that chain.  For each Hessian factor, its two distinct
creating primaries have different positions, so we can place it
strictly between its earliest and latest primary neighbors.  No edge
joins two Hessians, so these placements impose no competing condition.
Order Hessians occupying the same interval arbitrarily, and orient
every uncut edge from its earlier factor to its later factor.
The $b$ cut indices enter the first primary and leave the last one.

Every tensor has both an input and an output in this ordering.
An interior primary receives one original chain bundle and sends the
next; at the two endpoints the cut bundle supplies the missing side.
A Hessian receives an edge from an earlier neighbor and sends an edge
to a later one.  Thus all the matrix unfoldings used below are
nontrivial splits, exactly as required by the induction hypothesis.

Here is an explicit matrix-composition interpretation.  Immediately
before processing a factor, retain one copy of $\R^d$ for each index
whose earlier endpoint has been processed but whose later endpoint
has not.  The initial space is $(\R^d)^{\otimes b}$, belonging to
the cut inputs.  Permute these tensor coordinates so that the incoming
indices of the next factor are adjacent.  Apply that factor's matrix
unfolding to them and the identity to all other retained coordinates;
then permute the outputs into the new retained list.  Matrix
multiplication performs precisely the sums over the incoming indices.
Induction over the factors therefore accounts for every uncut index
once and leaves only the $b$ cut outputs at the end.  If $Q$ denotes
the resulting map from $(\R^d)^{\otimes b}$ to itself, then
\[
 \|Q\|_{\op}
 \le\prod_{i=1}^{2s}|\nabla^{j+e_i}V|_{\mathrm{split}}
      \prod_{v=1}^{n_H}|\nabla^{h_v}\nabla^2H_\rho|_{\mathrm{split}}.
\]
Indeed coordinate permutations and the intervening identity maps have
operator norm one.  Closing the original cut is now exactly a single
matrix trace, so the absolute value of this completely contracted term
is at most
\begin{equation}\label{eq:divergence-single-trace}
 |\tr Q|\le d^b\|Q\|_{\op}.
\end{equation}
This argument explains why other undirected cycles introduce no further
dimension factor: their coordinate sums have already been performed
by the matrix compositions.  Figure~\ref{fig:tensor-cut} shows a
possible ordering.

\begin{figure}[tb]
\centering
\begin{tikzpicture}[
  >=Latex,
  primary/.style={draw,rectangle,minimum width=9mm,minimum height=8mm},
  hessian/.style={draw,circle,minimum size=8mm},
  bundle/.style={->,line width=1.1pt},
  index/.style={->,line width=.5pt},
  font=\small]
 \node[primary] (p1) at (0,0) {$P_1$};
 \node[primary] (p2) at (2,0) {$P_2$};
 \node[hessian] (h) at (3,1.35) {$\mathcal H$};
 \node[primary] (p3) at (4,0) {$P_3$};
 \node[primary] (p4) at (6,0) {$P_4$};
 \coordinate (in) at (-1.2,0);
 \coordinate (out) at (7.2,0);
 \draw[bundle] (in) -- (p1);
 \draw[bundle] (p1) -- (p2);
 \draw[bundle] (p2) -- (p3);
 \draw[bundle] (p3) -- (p4);
 \draw[bundle] (p4) -- (out);
 \draw[index] (p1.north) to[out=80,in=180] (h.west);
 \draw[index] (h.east) to[out=0,in=95] (p3.north);
 \draw[index] (p2.south) to[out=-40,in=-140] (p4.south);
 \draw[index] (p2.south west) to[out=-70,in=-110] (p4.south east);
 \draw[dashed] (out) -- (7.2,-2.4) -- (-1.2,-2.4) -- (in);
 \node[fill=white,inner sep=3pt] at (3,-2.4)
   {the sole closing trace: cost $d^b$};
 \node[anchor=south] at (3,1.95)
   {order $P_1,P_2,\mathcal H,P_3,P_4$};
 \node[anchor=south] at (-.75,.60) {input};
 \node[anchor=south] at (6.75,.60) {output};
\end{tikzpicture}

\caption{A schematic term with four primary tensors
$P_i=\nabla^{j+e_i}V$ after removing the adjoint divergences.  Thick arrows are the surviving trace-chain
bundles; thin arrows are new contracted indices.  Here
$\mathcal H=\nabla^2H_\rho$ is the Hessian factor, inserted
between its first and last primary neighbors.  The arrows specify matrix inputs and outputs, so every factor
acts across a nontrivial split.  The dashed closure is the one trace
in \eqref{eq:divergence-single-trace}; its cost $d^b$ becomes bounded
after taking a sufficiently high moment root.}
\label{fig:tensor-cut}
\end{figure}

\paragraph{Uniform high-moment bookkeeping.}
The coordinate expansion and its contraction estimate are now complete.
It remains to verify that their constants allow $s$ to grow with $d$
and $j$.  Apply the induction hypothesis to every primary factor with
moment exponent $2s$, and use the pointwise Hessian estimates on all
other factors.  There are exactly $2s$ random primary factors, so H\"older's
inequality uses no larger exponent.  The power of $1/\sqrt a$ before
taking the $2s$-th root is at most
\[
 2s(k_V+j)+\sum_i e_i+\sum_v h_v
 \le 2s(k_V+j+1).
\]
Since $e_i,h_v\le2s$, all polynomial bases in the inductive estimates
are at most $C(s+j+\log(d+1))$.  The total exponent of those bases is
at most
\[
 C\left(2s(j+1)+\sum_i e_i+\sum_v(h_v+1)\right)
 \le C's(j+1),
\]
using \eqref{eq:network-degree-count}.  The term count $(Cs)^{Cs}$
is absorbed in the same bound.  The cut bundle uses only the original free slots of
$T=\nabla^*(\nabla^jV)$, not the new indices introduced by
integration by parts.  Consequently $b\le j+C_V$, where $C_V$
depends only on the fixed number of slots of $V$, and not on $s$.  We have proved
\[
 \E\|T\|_{\op}^{2s}
 \le d^{j+C_V}
       [C(s+j+\log(d+1))]^{Cs(j+1)}
       (\sqrt a)^{-2s(k_V+j+1)}.
\]
Choose $s$ to be the smallest positive integer satisfying $2s\ge p$
and $s\ge(j+C_V)\log(d+1)$, increasing $C_V$ if necessary.
Then $d^{(j+C_V)/(2s)}\le e^{1/2}$ and
\[
 s\le C_V\bigl(p+(j+1)\log(d+1)\bigr).
\]
The expression on the right is bounded by a fixed multiple of
$(p+j+\log(d+1))^2$.  Taking the $2s$-th root, using monotonicity
of probability-space $L^p$ norms, and increasing $C_E$ now gives
\eqref{eq:network} with $k=k_V+1$.

All integration-by-parts identities above are legitimate for each
finite $j,s,d$ and each $\rho<1$.  Indeed, the conditional-integral
regularity in Proposition~\ref{prop:gaussian-jets}, together with
the score formula, gives polynomial growth for every field and
derivative appearing in the finite expansion.  The laws $\xi_\rho$
have Gaussian tails.  Cutoffs therefore justify the identities with
vanishing boundary terms.  Uniformity in growing $s$ is supplied by
the explicit estimates after the identities, not by an unproved
uniform regularity assertion.  This completes the induction.
\end{proof}

\Needspace{9\baselineskip}
\subsection{One final vector dimension factor}

The preceding theorem controls every tensor expression while at least
two slots remain.  Only the last divergence producing a vector will
cost a $\sqrt d$ factor.  The following elementary identity makes
that cost explicit.

\begin{lemma}[Frobenius divergence estimate]
\label{lem:divergence-l2}
Let $\mu(dy)\propto e^{-H(y)}dy$, where $H\in C^2(\R^d)$
has bounded Hessian and $\mu$ has Gaussian tails.  Let $V$ be a smooth
tensor field with a distinguished divergence slot and with
$V,\nabla V\in L^2(\mu)$ in Frobenius norm.  Then
\begin{equation}
 \|\nabla^*V\|_{F,2}
 \le C\bigl(\|V\|_{F,2}+\|\nabla V\|_{F,2}\bigr),
 \label{eq:divergence-l2}
\end{equation}
where $C$ depends only on $\|\nabla^2H\|_{\op,\infty}$, and is
independent of $d$ and the number of remaining tensor slots.
For $H=H_\rho$ it is an absolute constant.
\end{lemma}

\begin{proof}
First take a compactly supported smooth vector field, so that the
output is scalar.  Integrating by parts twice gives
\[
 \int(\nabla^*V)^2\,d\mu
 =\int\sum_{i,j}
     \left[(\partial_iV_j)(\partial_jV_i)
                  +V_i(\partial_{ij}H)V_j\right]d\mu.
\]
The first term is bounded in absolute value by $|\nabla V|_F^2$,
and the second by $\|\nabla^2H\|_{\op}|V|^2$.  For tensor output,
apply the scalar identity to each fixed collection of remaining
indices and sum.  Taking a square root proves
\eqref{eq:divergence-l2} for compactly supported fields.  For general
$V$, multiply by smooth cutoffs $\chi_R$ equal to one on the ball of
radius $R$, with $|\nabla\chi_R|\le C/R$.  Then $\chi_RV\to V$
and $\nabla(\chi_RV)\to\nabla V$ in $L^2(\mu)$, while their
adjoint divergences converge pointwise.  Fatou's lemma proves the
stated estimate.  All fields to which we apply it have polynomially
growing derivatives and hence satisfy the required integrability.
\end{proof}

\begin{corollary}[Derivatives along transported paths]
\label{cor:path-derivatives}
For each fixed integer $m\ge1$,
\begin{equation}
 \|\mathcal D^mM_\rho\|_{L^2(\xi_\rho)}
 \le\polylog(d)L\sqrt d\,a^{-m}.
 \label{eq:path-mean}
\end{equation}
For fixed integers $q\ge1$, $m\ge0$, and fixed real
$2\le p<\infty$, every nontrivial split satisfies
\begin{equation}
 \bigl\|\,|\mathcal D^mU_q|_{\mathrm{split}}\,\bigr\|_{L^p(\xi_\rho)}
 \le\polylog(d)L(\sqrt a)^{-(q-1+2m)}.
 \label{eq:path-jet}
\end{equation}
The constants hidden in $\polylog(d)$ may depend on $m,q,p$
when these parameters occur, but are uniform in $\rho$, the primitive
center, scale, and oracle.
\end{corollary}

\begin{proof}
Iterate Lemma~\ref{lem:material-identities}, using the product rule
on contractions.  Replacing $U_q$ by $\nabla^*U_{q+1}$ raises weight
by two.  Replacing it by one of the contracted products in
\eqref{eq:material-jet-identity} does not raise weight: the two jet
weights sum to $q-1$.  The commutator
\eqref{eq:material-adjoint-identity} inserts $DM_\rho=\rho U_1$
of weight zero on the divergence slot.  Every additional jet factor
$L$ is accompanied by a factor $r$, so, after retaining one overall
$L$, the extra factors are powers of $l=rL\le1$.  Scalar
coefficients are fixed polynomials in $\rho$ and remain bounded
under the fixed number of differentiations.  Starting from a jet
$U_q$, all divergences retain at least two free slots, and all
contractions retain a slot of each factor.  The resulting terms are
therefore expressions of Theorem~\ref{thm:tensor-divergence}, of weight
at most $q-1+2m$.  Its case $j=0$ proves \eqref{eq:path-jet}.

Starting instead from $M$, the first identity is
$\mathcal DM=\nabla^*U_1$.  Its outermost divergence remains a
divergence to a vector under further material differentiations.
Inside it, the same construction gives finitely many normalized
expressions of weight at most $2(m-1)$, each with one overall factor
$L$.  Apply Lemma~\ref{lem:divergence-l2} at this outermost operation.
For a tensor with at least two slots, its Frobenius norm is at most
$\sqrt d$ times the operator norm of an unfolding that puts one
coordinate slot on one side: that matrix has at most $d$ singular
values.  Thus Theorem~\ref{thm:tensor-divergence}, for $j=0,1$ and
$p=2$, bounds the right side of \eqref{eq:divergence-l2} by
\[
 \polylog(d)L\sqrt d\,(\sqrt a)^{-(2m-1)}
 \le\polylog(d)L\sqrt d\,a^{-m}.
\]
This is \eqref{eq:path-mean}.
\end{proof}

To see the quantitative connection with quadrature, start the transport
at its correct marginal and put $Q_\rho=M_\rho(Y_\rho)$.
Proposition~\ref{prop:transport} ensures that $Y_\rho\sim\xi_\rho$,
so \eqref{eq:path-mean} bounds $\|Q_\rho^{(m)}\|_2$ for every
probability time.  On a cell $[\rho_0,\rho_1]\subset[0,1)$ of length
$h_c$, Taylor's integral remainder and the $L^2$ triangle inequality give,
for every $\rho\in[\rho_0,\rho_1]$,
\[
 \left\|Q_\rho-
       \sum_{h=0}^{m-1}\frac{(\rho-\rho_0)^h}{h!}Q_{\rho_0}^{(h)}\right\|_2
 \le \frac{h_c^m}{m!}\sup_{t\in[\rho_0,\rho_1]}\|Q_t^{(m)}\|_2
 \le C_m\polylog(d)L\sqrt d\,
             h_c^m(1-\rho_1^2)^{-m}.
\]
In particular, for a mesh ratio $\theta>0$, the condition
$h_c\le\theta(1-\rho_1)$ makes this remainder at most
$C_m\polylog(d)L\sqrt d\,\theta^m$.
This is a supremum of deterministic RMS bounds; it does not require
an expected supremum over a random trajectory.  The proof of
Proposition~\ref{prop:probability-circuit} uses this cancellation
between cell length and endpoint singularity in polynomial quadrature.
Before constructing those finite formulas, we use the material identity
itself to produce a stationary flow whose path integral represents a
noisy mean.
